一个有效的两模有增益--损耗量子系统（或光学两通道耦合模系统）在 \(\hbar=1\) 单位下由非厄米矩阵
\[
H=\begin{pmatrix}i\gamma&J\\ J&-i\gamma\end{pmatrix}
=J\sigma_x+i\gamma\sigma_z,\qquad J>0,\quad\gamma\ge0
\]
支配，演化满足 \(i\partial_t|\psi\rangle=H|\psi\rangle\)。定义空间交换算符 \(P=\sigma_x\) 和反线性时间反演 \(T\) 为复共轭。初态取 \(|\psi(0)\rangle=(1,0)^T\)，普通范数记作 \(N(t)=\psi^\dagger\psi\)。注意这个有效非厄米演化不应被当作闭合量子系统的幺正传播。
(1) 证明 \(PT\) 对称，严格计算本征值、实谱区与虚谱区；在 \(\gamma=J\) 证明代数重数为2但几何重数为1，构造其二阶例外点的广义本征结构。
(2) 在三个参数区（\(\gamma<J\)、\(\gamma=J\)、\(\gamma>J\)）分别求完整时间演化算符和给定初态的两分量与普通范数，说明为何标准范数不守恒以及例外点为什么出现多项式而非通常指数传播。
(3) 在 \(0\le\gamma<J\) 构造一个显式正定厄米度量矩阵 \(\eta\ne I\)，证明 \(H^\dagger\eta=\eta H\) 并得到严格守恒的加权范数；计算度量的两个特征值并解释为什么它在例外点不再严格正定。
(4) 固定 \(J\)，研究从实谱一侧接近例外点时 \(\gamma=J-\varepsilon\)、\(\varepsilon\downarrow0\) 的能级劈裂主项与导数灵敏度；对破缺区求长时间普通范数增长指数，并说明实验上有效增益/损耗导致的非幺正传播与真正量子态条件概率归一化的区别。